\documentclass[12pt]{article}
\usepackage[utf8]{inputenc}
\usepackage{amsmath}
\usepackage{amssymb}
\usepackage{graphicx}
\usepackage{booktabs}
\usepackage{listings}
\usepackage{xcolor}
\usepackage{hyperref}
\usepackage{geometry}
\geometry{margin=1in}

\lstset{
    language=Python,
    basicstyle=\ttfamily\small,
    keywordstyle=\color{blue}\bfseries,
    commentstyle=\color{green!60!black},
    stringstyle=\color{red},
    numbers=left,
    numberstyle=\tiny\color{gray},
    stepnumber=1,
    numbersep=5pt,
    backgroundcolor=\color{gray!10},
    frame=single,
    breaklines=true,
    breakatwhitespace=false,
    tabsize=4,
    showstringspaces=false
}

\title{Line Loss Costs/Benefits Calculation Documentation}
\author{}
\date{}

\begin{document}

\maketitle

\section{Introduction}

The \texttt{line\_loss\_costs.py} script calculates line loss costs for greenfield projects and line loss benefits for reconductoring projects. For greenfield projects, line losses represent a cost (energy that must be generated to compensate). For reconductoring projects, reduced losses compared to the old conductor represent a benefit.

\section{Method Overview}

\subsection{Purpose}

This script calculates line loss costs/benefits by:
\begin{enumerate}
    \item Determining if project is greenfield or reconductoring
    \item For greenfield: Calculating losses and converting to costs
    \item For reconductoring: Comparing old vs. new conductor losses to calculate benefits
    \item Computing present values using appropriate discount rates
\end{enumerate}

\subsection{Calculation Approach}

\textbf{Greenfield Projects:}
\begin{align}
\text{Annual Loss Cost} &= \text{Losses (MWh/year)} \times \text{Electricity Price (\$/MWh)} \\
\text{Lifetime Nominal} &= \text{Annual Cost} \times \text{Project Lifetime} \\
\text{Present Value} &= \sum_{t=1}^{N} \frac{\text{Annual Cost}}{(1 + r)^{t + \text{start\_year} - 1}}
\end{align}

\textbf{Reconductoring Projects:}
\begin{align}
\text{Loss Reduction} &= \text{Old Losses} - \text{New Losses} \\
\text{Annual Benefit} &= \text{Loss Reduction} \times \text{Electricity Price} \\
\text{Lifetime Benefit} &= \text{Annual Benefit} \times \text{Project Lifetime} \\
\text{Present Value} &= \sum_{t=1}^{N} \frac{\text{Annual Benefit}}{(1 + r)^{t + \text{start\_year} - 1}}
\end{align}

The script uses a normalized comparison method that accounts for capacity differences between old and new configurations.

\subsection{Inputs and Outputs}

\textbf{Inputs:}
\begin{itemize}
    \item Project technical details (construction type, AC/DC, capacity, conductor type, reconductoring status)
    \item Old configuration details (for reconductoring: old capacity, old conductor type, old AC/DC)
    \item Physical details (line length)
    \item Circuit and resistance details
    \item Electricity price
    \item Social discount rate
    \item Project lifetime, delay years, construction years
\end{itemize}

\textbf{Outputs:}
\begin{itemize}
    \item Annual cost/benefit
    \item Lifetime nominal cost/benefit
    \item Present value cost/benefit
\end{itemize}

\subsection{Integration with Other Scripts}

This script integrates with:
\begin{itemize}
    \item \texttt{energy\_losses.py} --- \texttt{get\_total\_energy\_losses} (primary config), \texttt{load\_physical\_details}, \texttt{load\_circuit\_and\_resistance\_details}, \texttt{calculate\_phase\_current}, \texttt{full\_load\_adjusted}, \texttt{calculate\_line\_losses}
    \item \texttt{smart\_loaders.py} --- \texttt{load\_financing\_social\_discount\_rate}, \texttt{load\_project\_technical\_details}, \texttt{get\_project\_data\_raw}
    \item \texttt{financial\_utils.py} --- \texttt{calculate\_present\_value}, \texttt{calculate\_cod\_year}
    \item \texttt{calculation\_utils.py} --- \texttt{to\_percent}, \texttt{from\_percent}, \texttt{build\_category\_string}, \texttt{calculate\_converter\_losses}
    \item \texttt{constants.py} --- \texttt{HOURS\_PER\_YEAR}
    \item \texttt{smart\_output.py} --- \texttt{CTCCOutputManager} (\texttt{add\_line\_loss\_costs}, \texttt{write\_batch\_summary})
\end{itemize}

\section{Key Functions}

\subsection{\texttt{load\_project\_details()}}

Loads project technical details with line\_loss\_costs specific fields.

\textbf{Returns:}
\begin{itemize}
    \item Multiple values including construction type, AC/DC, capacity, conductor type, utilization, reconductoring status, timeline parameters
\end{itemize}

\subsection{\texttt{calculate\_configuration\_losses(...)}}

Calculates line losses (conductor only) for a given configuration. Used for reconductoring baseline (old) and for greenfield comparison configurations. Optional \texttt{voltage\_kv\_override} allows using a different voltage (e.g., reconductoring: old tower voltage).

\textbf{Parameters:}
\begin{itemize}
    \item \texttt{construction\_type}, \texttt{ac\_dc}, \texttt{capacity\_mw}, \texttt{conductor\_type}, \texttt{converter\_type}, \texttt{line\_utilization\_percent}, \texttt{project\_lifetime}
    \item \texttt{voltage\_kv\_override} (optional): If set, used instead of category lookup voltage
\end{itemize}

\textbf{Returns:}
\begin{itemize}
    \item \texttt{losses\_mwh\_per\_year}, \texttt{lifetime\_losses\_mwh} (conductor losses only; converter losses computed separately when needed)
\end{itemize}

\textbf{Key Logic:}
\begin{lstlisting}
category = build_category_string(construction_type, ac_dc, capacity_mw, conductor_type, converter_type)
line_length = load_physical_details()
circuit = load_circuit_and_resistance_details(category)
voltage_kv = voltage_kv_override if voltage_kv_override is not None else circuit.voltage_kv
capacity_mw_str = f"{capacity_mw}MW"
phase_current = calculate_phase_current(capacity_mw_str, voltage_kv, circuit.number_of_phases, circuit.number_of_circuits_poles, ac_dc)
full_load_adj = full_load_adjusted(line_utilization_percent)
losses_mwh_per_year, lifetime_losses_mwh, ... = calculate_line_losses(phase_current, full_load_adj, circuit.AC_75_resistance, circuit.DC_20_resistance, ac_dc, ...)
\end{lstlisting}

\subsection{\texttt{main()}}

Main execution function that orchestrates the line loss cost/benefit calculation.

\textbf{Workflow for Greenfield:}
\begin{enumerate}
    \item Gets losses via \texttt{get\_total\_energy\_losses()} (line + converter)
    \item Converts to costs (losses × \texttt{baseline\_electricity\_price}), \texttt{start\_year} = \texttt{calculate\_cod\_year(delay\_years, construction\_years)}
    \item Calculates present value; writes results (including optional line/converter breakdown) to CSV
\end{enumerate}

\textbf{Workflow for Reconductoring:}
\begin{enumerate}
    \item Baseline (old) losses: \texttt{calculate\_configuration\_losses(old\_...)} + \texttt{calculate\_converter\_losses} for old config; new (primary) losses: \texttt{get\_total\_energy\_losses()}
    \item Counterfactual: old conductor at new capacity via \texttt{calculate\_configuration\_losses} with new capacity
    \item Benefit = (baseline or counterfactual total losses − new total losses) × price; PV and CSV as benefits
\end{enumerate}

\section{Example}

The following example demonstrates a typical usage:

\begin{lstlisting}
# Example 1: Greenfield Project
# 500 MW Overhead AC line
# 
# Inputs:
#   - Capacity: 500 MW
#   - Line utilization: 70%
#   - Line length: 100 miles
#   - Losses: 5,000 MWh/year
#   - Electricity price: $50/MWh
#   - Project lifetime: 50 years
#   - Start year: 3 (after 2-year construction)
#
# Calculation:
#   Annual cost = 5,000 * $50 = $250,000/year
#   Lifetime nominal = $250,000 * 50 = $12.5M
#   PV = $250,000 * PV annuity (50 years, 3%, start year 3) = $8.2M
#
# Output:
#   - Annual cost: $250,000
#   - Lifetime nominal: $12.5M
#   - Present value: $8.2M

# Example 2: Reconductoring Project
# Upgrade 200 MW line to 350 MW
# 
# Inputs:
#   - Old: 200 MW, Standard Aluminum, losses: 3,000 MWh/year
#   - New: 350 MW, Advanced Aluminum, losses: 4,000 MWh/year
#   - Counterfactual (old conductor @ 350 MW): 5,500 MWh/year
#   - Electricity price: $50/MWh
#
# Calculation (Normalized Method):
#   Loss reduction = 5,500 - 4,000 = 1,500 MWh/year
#   Annual benefit = 1,500 * $50 = $75,000/year
#   Lifetime benefit = $75,000 * 50 = $3.75M
#   PV = $75,000 * PV annuity (50 years, 3%) = $1.8M
#
# Output:
#   - Annual benefit: $75,000
#   - Lifetime benefit: $3.75M
#   - Present value: $1.8M
\end{lstlisting}

\textbf{Integration Example:}

\begin{lstlisting}
# When called from ctcc.py (main() builds results):
csv_manager = CTCCOutputManager()
results = {
    "annual_cost": total_annual_cost,
    "total_nominal": total_lifetime_nominal,
    "total_afudc": 0,
    "total_pv": total_pv_loss_cost,
    "line_loss_mwh_yr": losses_mwh_per_year,
    "converter_loss_mwh_yr": total_converter_losses_mwh,
    "total_loss_mwh_yr": total_losses_mwh_per_year,
    "line_annual_cost": ..., "converter_annual_cost": ...,
    "line_cost_pv": ..., "converter_cost_pv": ...,
}
csv_manager.add_line_loss_costs(results)
csv_manager.write_batch_summary()
\end{lstlisting}

\section{Key Implementation Details}

\subsection{Greenfield vs. Reconductoring}

The script handles two distinct cases:
\begin{itemize}
    \item \textbf{Greenfield}: Line losses are a cost (energy that must be generated)
    \item \textbf{Reconductoring}: Reduced losses are a benefit (compared to old conductor)
\end{itemize}

\subsection{Comparison Methods for Reconductoring}

The script implements three comparison methods:
\begin{enumerate}
    \item \textbf{Direct Comparison}: Old losses @ old capacity vs. new losses @ new capacity
    \item \textbf{Counterfactual Comparison}: Old conductor @ new capacity vs. new conductor @ new capacity (used for output)
    \item \textbf{Normalized Comparison}: Loss percentages weighted by delivered energy
\end{enumerate}

The normalized method is used for the final output as it accounts for capacity differences.

\subsection{Counterfactual Baseline}

For reconductoring, the script calculates what losses would be if the old conductor were used at the new capacity. This provides a fair comparison by isolating the effect of the conductor upgrade.

\subsection{Present Value Timing}

Costs/benefits start when the project becomes operational (after delay and construction periods). The present value calculation discounts annual payments back to the base year.

\subsection{AFUDC Treatment}

Line loss costs/benefits are not AFUDC-eligible because they are operational costs/benefits, not construction costs.

\end{document}

